\section{Retaining the analytic tail instead of discarding it (NS-82)}
\label{sec:ns82-tail-compression}

The large sufficient cutoff in NS-81 is not a necessary cost. Its leading
logarithmic tail can instead be retained exactly as a rank-two form. The
following construction uses physical cell integrals, not the truncated
sampling matrix of NS-81. It still gives no cofinal error estimate.

For $f_c=\sum_{n\leq N}c_n\sigma_n$, set
\[
 u=\tfrac12\sum c_n,\qquad
 v=\tfrac12\sum c_n\log(2\pi/n),\qquad
 S_2(N)=\sum_{n\leq N}n^2=\frac{N(N+1)(2N+1)}6.
\]
Keep the harmonic number $H_N$ and constant $C_N=48H_NN^3(N+2)$
from NS-81. For an integer $M\geq N+1$, define the finite quadratic form
\begin{equation}
 T_M(c)=\int_{1/M}^{\infty}|f_c(x)|^2\,dx
       +\frac{q^2+2uq+2u^2}{M},\qquad q=u\log M+v.
 \label{eq:ns82-comparator}
\end{equation}
The first term consists of the exact exterior contribution and cells
$1\leq m<M$; the second is the integral of the logarithmic main tail.
Thus the added tail has rank at most two in coefficient space. Exact
cell formulas are those of NS-78, with $\alpha=0$ and the sign of the
residual reversed; their squares are unchanged.

\begin{proposition}[Complete remainder control for a finite physical comparator]
\label{prop:ns82-tail-compression}
Put $\delta^2=4C_NS_2(N)/(27M^3)$. If $\delta<1$, then
\begin{equation}
 (1-\delta)^2G_N\preceq T_M\preceq(1+\delta)^2G_N.
 \label{eq:ns82-comparison}
\end{equation}
In particular the explicit integer choice
\[
 M_N=\max\left\{N+1,
 \left\lceil\left(\frac{16C_NS_2(N)}{27}\right)^{1/3}\right\rceil\right\}
 =O\bigl(N^{7/3}\log^{1/3}(2N)\bigr)
\]
satisfies $\delta\leq1/2$ and gives
$G_N/4\preceq T_{M_N}\preceq9G_N/4$.
The same constants survive old-space minimization, so the projected
condition ratio is at most nine. An old $N$-to-$2N$ block uses $M_{2N}$.
\end{proposition}
\begin{proof}
NS-78's complete Bernoulli remainder gives, for $t\geq M\geq N$,
\[
 f_c(1/t)=u\log t+v+e_c(t),\qquad
 |e_c(t)|\leq\frac{\sum n|c_n|}{6t}
            \leq\frac{\sqrt{S_2(N)}\,\|c\|_2}{6t}.
\]
Consequently the norm of the full discarded remainder is bounded by
\[
 \int_M^\infty\frac{|e_c(t)|^2}{t^2}\,dt
 \leq\frac{S_2(N)}{108M^3}\|c\|_2^2
 \leq\frac{4C_NS_2(N)}{27M^3}\|f_c\|_2^2,
\]
where the final step uses NS-81's ordinary Gram floor. Replace $f_c$
only on $0<x<1/M$ by $u\log(1/x)+v$. The resulting function has squared
norm exactly $T_M(c)$. The triangle and reverse-triangle inequalities
therefore give~\eqref{eq:ns82-comparison}, retaining the cross term
between the main tail and its remainder. The formula for $M_N$ forces
$\delta^2\leq1/4$. Taking the infimum over old coordinates preserves
both quadratic-form constants, as in NS-80. No equality of old
minimizers is required.
\end{proof}

For the true projected correlations, the NS-80 exact line-search argument
now gives one-step capture at least $9/25$ of full available gain; ten
ideal steps capture at least
\[
 1-(4/5)^{20}>.9884.
\]
The comparator uses finitely many elementary cell integrals and a closed
rank-two tail; no infinite comparator or unevaluated tail is hidden.
The cutoff is calculated using exact rational harmonic numbers and an
integer cube-root inequality. Two-precision Arb checks additionally
verify the complete remainder constants at $N=16,64,256,512$.
These are scalar cutoff certificates only: the comparator matrices and
preconditioned iterations are not assembled or executed here.

This improves the sufficient construction cost, not the approximation
rate or the original signed Weil floor. The new size comparison is
uniform over all vectors and projected blocks; its useful descent is
still measured against available block gain, which may be smaller than
the entire residual error. Cofinal arithmetic decay, G2 and RH remain open.
